% ---------------------------------------------------------------------
% Argomento: rete mista con E e I note e una resistenza del parallelo
%            incognita (R3): si risale da Req a R23, poi alle correnti, e KVL
% Livello: intermedio-avanzato   |   Tempo stimato: 20 min
% Lezioni: 7/10 (serie, parallelo, partitori), 9/10 (KCL e KVL)
% ---------------------------------------------------------------------
\begin{esercizio}
\dato{E}{230}{230}{230}           % generatore [V]
\dato{I}{5.4}{5.3}{5.2}           % corrente erogata dal generatore [A]
\dato{R1}{12}{11}{7}              % [ohm]
\dato{R2}{48}{88}{66}
\dato{R4}{14}{5}{6}
\calcola[R_{eq}]{Req}{\x{E}/\x{I}}
\calcola[R_{23}]{R23}{\x{Req}-\x{R1}-\x{R4}}
\calcola[V_{23}]{V23}{\x{R23}*\x{I}}
\calcola[I_2]{I2}{\x{V23}/\x{R2}}
\calcola[I_3]{I3}{\x{I}-\x{I2}}
\calcola[R_3]{R3}{\x{V23}/\x{I3}}
\calcola[V_1]{V1}{\x{R1}*\x{I}}
\calcola[V_4]{V4}{\x{R4}*\x{I}}

Il generatore eroga la corrente $I$. Sono note $R_1$, $R_2$, $R_4$ ma non $R_3$.

\begin{center}
\begin{minipage}{0.52\linewidth}\centering
  \begin{circuitikz}[x=0.85cm, y=0.5cm]
    \ctikzset{bipoles/length=1cm}
    \draw (0,0) to[vsource, l=$E$, n=G] (0,3)
          to[R, l=$R_1$, i=$I$] (3,3) node[circ]{} -- (5.5,3)
          (3,3) to[R, l_=$R_2$] (3,0) node[circ]{}
          (5.5,3) to[R, l=$R_3$] (5.5,0) -- (3,0)
          (3,0) to[R, l=$R_4$] (0,0);
    \morsetti{G}
  \end{circuitikz}
\end{minipage}\hfill
\begin{minipage}{0.4\linewidth}
  \renewcommand{\arraystretch}{1.1}
  \begin{tabular}{r@{${}={}$}l}
    $E$   & \mis{E}{\volt}\\
    $I$   & \mis{I}{\ampere}\\
    $R_1$ & \mis{R1}{\ohm}\\
    $R_2$ & \mis{R2}{\ohm}\\
    $R_4$ & \mis{R4}{\ohm}
  \end{tabular}
\end{minipage}
\end{center}

\vspace{-6pt}
\begin{parts}
  \part[1] Calcola la resistenza equivalente $R_{eq}$ vista dal generatore.
  \part[2] Calcola la resistenza $R_{23}$ del parallelo tra $R_2$ e $R_3$.
  \part[1] Calcola la tensione $V_{23}$ ai capi del parallelo.
  \part[1] Calcola la corrente $I_2$ in $R_2$.
  \part[1] Calcola la corrente $I_3$ in $R_3$.
  \part[2] Calcola la resistenza $R_3$.
  \part[2] Calcola le tensioni $V_1$ e $V_4$ ai capi di $R_1$ e $R_4$.
  \part[2] Applica la seconda legge di Kirchhoff alla maglia $E$, $R_1$, $R_2$, $R_4$
    e verifica che i conti tornino.
\end{parts}

\begin{solution}
  $R_{eq} = E/I \approx \mis{Req}{\ohm}$\\[4pt]
  $R_{eq} = R_1 + R_{23} + R_4 \;\Rightarrow\; R_{23} = R_{eq} - R_1 - R_4 \approx \mis{R23}{\ohm}$\\[4pt]
  $V_{23} = R_{23}\,I \approx \mis{V23}{\volt}$\\[4pt]
  $I_2 = V_{23}/R_2 \approx \mis{I2}{\ampere}$\\[4pt]
  KCL: $I_3 = I - I_2 \approx \mis{I3}{\ampere}$\\[4pt]
  $R_3 = V_{23}/I_3 \approx \mis{R3}{\ohm}$\\[4pt]
  $V_1 = R_1\,I \approx \mis{V1}{\volt}$,\qquad $V_4 = R_4\,I \approx \mis{V4}{\volt}$\\[4pt]
  KVL: $E = V_1 + V_{23} + V_4 = \mis{V1}{\volt} + \mis{V23}{\volt} + \mis{V4}{\volt} \approx \mis{E}{\volt}$ \quad (i conti tornano)
  \criterio{a) 1 pt; b) 1 pt serie $R_{eq}=R_1+R_{23}+R_4$, 1 pt risultato; c) 1 pt;
    d) 1 pt; e) 1 pt KCL; f) 1 pt formula, 1 pt risultato; g) 1 pt per ciascuna tensione; h) 1 pt equazione $E=V_1+V_{23}+V_4$, 1 pt verifica numerica. In alternativa si può usare
    $R_3 = R_2R_{23}/(R_2-R_{23})$: stessi punti. Unità mancanti: $-0{,}5$ pt una volta sola.}
\end{solution}
\end{esercizio}
